\section{Описание}

Требуется реализовать поиск одного или нескольких образцов в тексте с использованием алгоритма \textbf{Кнута-Морриса-Пратта} (КМП) с построением функции переходов через Z-функцию по методу Гусфилда.

Основная идея алгоритма КМП заключается в том, что при каждом несовпадении символа не требуется начинать поиск заново. Вместо этого используется функция неудач $F$, которая позволяет совершать максимально длинные сдвиги паттерна, используя информацию о его строении.

Для эффективной работы с потоками данных большого объёма используется \textbf{потоковая обработка}: текст читается построчно, разбивается на слова, и каждое слово обрабатывается без накопления всех данных в памяти.

Поскольку паттерн состоит из нескольких слов, а вхождение может начинаться с последнего прочитанного слова буфера, применяется \textbf{кольцевой буфер} размером $n$ (количество слов в паттерне) для сохранения метаданных о позициях последних обработанных слов.

Для оптимизации вычисления функции неудач используется алгоритм \textbf{Гусфилда}, который на основе $Z$-функции быстро строит таблицу сдвигов за линейное время $O(n)$.

\pagebreak

\section{Исходный код}

Реализация состоит из нескольких ключевых компонентов, каждый отвечающий за отдельный аспект алгоритма.

\subsection{Структура Pos}

Вспомогательная структура для хранения метаданных о позиции слова в тексте. Содержит:
\begin{itemize}
    \item \texttt{size\_t line} — номер строки;
    \item \texttt{size\_t pos} — номер слова в строке.
\end{itemize}

\subsection{Класс FArray}

Класс, отвечающий за предварительную обработку паттерна. Реализует построение массива сдвигов "функцию неудач" $F$. Данный подход позволяет за линейное время $O(n)$ вычислить все необходимые переходы для автомата КМП.

Вычисление функции неудач разделяется на 3 основных этапа:
\begin{enumerate}
\item Построение Z-блоков (Z-функции): На этом этапе для каждой позиции паттерна вычисляется длина наибольшего общего префикса между всем паттерном и его суффиксом, начинающимся с этой позиции.
\item Построение функции $sp'$: Данная функция определяет длину максимального строгого префикса паттерна, который одновременно является и его суффиксом для текущей подстроки, при условии, что символы, следующие за этим префиксом и суффиксом, не совпадают.
\item Построение функции неудач $F$: Финальный этап, на котором на основе полученных значений $sp'$ формируется итоговый массив сдвигов. Функция $F$ указывает, к какому состоянию (или какому символу паттерна) необходимо перейти в случае возникновения несоответствия при поиске.
\end{enumerate}
\subsection{Потоковая обработка и поиск}

Основная логика сосредоточена в функции \texttt{main} и локальной лямбда-функции \texttt{process word}.

Программа работает в следующем порядке:

\begin{enumerate}
    \item \textbf{Чтение и парсинг паттерна:} Первая строка нормализуется (нижний регистр) и разбивается на слова.
    \item \textbf{Инициализация автомата:} Строится массив $F$ для вычисления переходов.
    \item \textbf{Процесс поиска:} Каждое слово текста сопоставляется с текущим состоянием автомата.
    \item \textbf{Обновление состояния:} Переменная $p$ отслеживает количество совпавших слов. При несовпадении используется функция $F$.
    \item \textbf{Вывод результатов:} При полном совпадении ($p == n + 1$) координаты извлекаются из кольцевого буфера.
\end{enumerate}

\subsection{Кольцевой буфер}

Кольцевой буфер хранит координаты последних $n$ обработанных слов. Это необходимо для корректного вывода позиции начала найденного паттерна в условиях потокового чтения.

Индекс в кольцевом буфере циклически увеличивается: \texttt{ring\_idx = (ring\_idx + 1) \% n}. Когда найдено совпадение, начало паттерна находится в позиции \texttt{(ring\_idx + 1) \% n}.

\begin{table}[!ht]
\centering
\begin{tabular}{|p{7.5cm}|p{7.5cm}|}
\hline
\rowcolor{lightgray}
\multicolumn{2}{|c|}{main.cpp} \\
\hline
\texttt{FArray(const std::vector<std::string>\& pattern)} & Конструктор: вычисление Z-функции и построение таблицы $F$ \\
\hline
\rowcolor{lightgray}
\multicolumn{2}{|c|}{main.cpp} \\
\hline
\texttt{void process\_word(size\_t l, size\_t w)} & Лямбда-функция обработки одного слова и обновления состояния автомата КМП \\
\hline
\rowcolor{lightgray}
\multicolumn{2}{|c|}{main.cpp} \\
\hline
\texttt{std::vector<Pos> ring\_buffer} & Кольцевой буфер для хранения координат последних $n$ слов \\
\hline
\end{tabular}
\end{table}

\pagebreak

\begin{lstlisting}[language=C++]
#include <iostream>
#include <string>
#include <vector>
#include <cstdint>
#include <algorithm>

struct Pos {
    size_t line;
    size_t pos;
};

class FArray {
    std::vector<size_t> F;

public:
    explicit FArray(const std::vector<std::string>& pattern) {
        size_t n = pattern.size();
        if (n == 0) return;

        std::vector<size_t> Z(n, 0);
        size_t L = 0, R = 0;
        for (size_t i = 1; i < n; ++i) {
            if (i <= R) {
                Z[i] = std::min(R - i + 1, Z[i - L]);
            }
            while (i + Z[i] < n && pattern[Z[i]] == pattern[i + Z[i]]) {
                ++Z[i];
            }
            if (i + Z[i] - 1 > R) {
                L = i;
                R = i + Z[i] - 1;
            }
        }
        std::vector<size_t> sp_prime(n, 0);
        for (int64_t j = static_cast<int64_t>(n) - 1; j >= 1; --j) {
            if (Z[j] > 0) {
                size_t i = j + Z[j] - 1;
                sp_prime[i] = Z[j];
            }
        }
        F.resize(n + 2, 0);
        F[1] = 1;
        for (size_t i = 2; i <= n + 1; ++i) {
            F[i] = sp_prime[i - 2] + 1;
        }
    }

    inline size_t operator[](size_t p) const {
        return F[p];
    }
};

int main() {
    std::ios_base::sync_with_stdio(false);
    std::cin.tie(NULL);

    std::string pattern_line;
    if (!std::getline(std::cin, pattern_line) || pattern_line.empty()) {
        return 0;
    }
    std::vector<std::string> pattern;
    std::string temp;
    for (char ch : pattern_line) {
        if (std::isspace(ch)) {
            if (!temp.empty()) {
                for (char &c : temp) c = std::tolower(c);
                pattern.push_back(temp);
                temp.clear();
            }
        } else {
            temp += ch;
        }
    }
    if (!temp.empty()) {
        for (char &c : temp) c = std::tolower(c);
        pattern.push_back(temp);
    }
    size_t n = pattern.size();
    if (n == 0) return 0;

    FArray F(pattern);

    std::vector<Pos> ring_buffer(n);

    size_t ring_idx = 0;
    size_t line_num = 1;

    size_t p = 1;

    std::string current_word;
    current_word.reserve(20); 

    auto process_word = [&](size_t l, size_t w) {
        if (current_word.empty()) return;

        for (char &c : current_word) c = std::tolower(c);

        ring_buffer[ring_idx] = {l, w};

        while (true) {
            if (p <= n && current_word == pattern[p - 1]) {
                p++;
                break;
            }
            if (p == 1) break;
            p = F[p];
        }
        if (p == n + 1) {
            size_t match_start_idx = (ring_idx + 1) % n;
            std::cout << ring_buffer[match_start_idx].line << ", " << ring_buffer[match_start_idx].pos << "\n";
            p = F[p];
        }
        ring_idx = (ring_idx + 1) % n;
        current_word.clear();
    };

    std::string line;
    while (std::getline(std::cin, line)) {
        size_t word_in_line = 1;

        for (size_t i = 0; i < line.length(); ++i) {

            char ch = line[i];

            if (std::isspace(ch)) {
                if (!current_word.empty()) {
                    process_word(line_num, word_in_line);
                    word_in_line++;
                }
            } else {
                current_word += ch;
            }
        }
        if (!current_word.empty()) {
            process_word(line_num, word_in_line);
        }
        line_num++;
    }
    return 0;
}
\end{lstlisting}

\pagebreak

\subsection{Анализ сложности}

\begin{itemize}
    \item \textbf{Временная сложность предварительной обработки:} $O(n)$, где $n$ — количество слов в паттерне.
    \item \textbf{Временная сложность поиска:} $O(m + n)$, где $m$ — общее количество слов в тексте. Каждое слово текста обрабатывается ровно один раз.
    \item \textbf{Пространственная сложность:} $O(n)$ для хранения паттерна, функции $F$ и кольцевого буфера.
\end{itemize}

Благодаря кольцевому буферу размером $n$ и потоковой обработке программа способна работать с текстами произвольного размера, не накапливая их целиком в памяти.

\pagebreak

\section{Консоль}

\begin{alltt}
root@workspaces:/workspaces/Discrete-Analysis-MAI# cmake -B build
root@workspaces:/workspaces/Discrete-Analysis-MAI# cmake --build build
root@workspaces:/workspaces/Discrete-Analysis-MAI# ./build/lab4/lab4

cat dog cat dog bird
CAT dog CaT Dog Cat DOG bird CAT
dog cat dog bird
1, 3
1, 8

\end{alltt}

\pagebreak